\documentclass[10pt]{article}
\usepackage[T1]{fontenc}
\usepackage[utf8]{inputenc}
\usepackage[french]{babel}
\usepackage{amsmath,amssymb,mathrsfs,mathpazo}
\usepackage{geometry}
\geometry{a4paper,margin=2cm}
\usepackage{graphicx,array,tabularx,tabularray,colortbl,multirow,multicol,booktabs,enumitem,pifont,ulem,fancybox,diagbox,cancel}
\usepackage{pstricks,pstricks-add,pst-plot,pst-tree,pst-node,pst-text,pst-func,pst-math,pst-bspline,pst-3dplot,pst-eucl,pst-solides3d}
\let\pscalculate\undefined
\usepackage{pst-bezier}
\usepackage{tikz,pgfplots,tkz-tab}
\usetikzlibrary{arrows,arrows.meta,patterns,calc,positioning,trees,decorations.markings,shapes.geometric,shapes.misc,shapes.arrows,intersections,decorations.pathmorphing,decorations.pathreplacing,backgrounds,shadows,babel}
\pgfplotsset{compat=1.18}
\usepgfplotslibrary{fillbetween}
\usepackage[np]{numprint}
\usepackage{listings,makecell,pas-tableur,twoopt}
\let\barre\undefined
\usepackage[pstricks]{tablvar}
\definecolor{darkgreen}{RGB}{0,100,0}
\providecommand{\red}{\color{red}}
\providecommand{\blue}{\color{blue}}
\providecommand{\green}{\color{green}}
\providecommand{\black}{\color{black}}
\providecommand{\white}{\color{white}}
\providecommand{\gray}{\color{gray}}
\providecommand{\cyan}{\color{cyan}}
\providecommand{\magenta}{\color{magenta}}
\providecommand{\yellow}{\color{yellow}}
\providecommand{\violet}{\color{violet}}
\providecommand{\orange}{\color{orange}}
\providecommand{\euro}{\texteuro}
\providecommand{\EUR}{\texteuro}
\providecommand{\ieme}{\textsuperscript{e}}
\providecommand{\degres}{\ensuremath{^\circ}}
\providecommand{\degree}{\ensuremath{^\circ}}
\providecommand{\vv}[1]{\overrightarrow{#1}}
\providecommand{\Coord}[1]{#1}
\providecommand{\decosix}{}
\providecommand{\starredbullet}{\ensuremath{\bullet}}
\pagestyle{empty}
\setlength{\parindent}{0pt}
\usepackage{adjustbox}\DeclareUnicodeCharacter{2236}{:}\DeclareUnicodeCharacter{2190}{\ensuremath{\leftarrow}}
\providecommand{\R}{\mathbb{R}}
\providecommand{\N}{\mathbb{N}}
\providecommand{\Z}{\mathbb{Z}}
\providecommand{\Q}{\mathbb{Q}}
\providecommand{\D}{\mathbb{D}}
\providecommand{\C}{\mathbb{C}}
\providecommand{\e}{{\mathrm{e}}}
\providecommand{\vect}[1]{\overrightarrow{#1}}
\providecommand{\vectt}[1]{\overrightarrow{\mathrm{#1}}}
\providecommand{\vvt}[1]{\overrightarrow{\mathrm{#1}}}
\providecommand{\np}[2][]{{#2~#1}}
\providecommand{\vv}[1]{{\overrightarrow{#1}}}
\renewcommand{\d}{{\mathrm{d}}}
\providecommand{\degres}{{^\circ}}
\providecommand{\degre}{{^\circ}}
\providecommand{\blue}{\color{blue}}
\providecommand{\red}{\color{red}}
\providecommand{\enskip}{\quad}
\providecommand{\llbracket}{\lbrack\!\lbrack}
\providecommand{\rrbracket}{\rbrack\!\rbrack}
\providecommand{\euro}{{\text{€}}}
\providecommand{\no}{{\text{n°}}}
\providecommand{\interval}[2]{{[#1;#2]}}
\providecommand{\intervFO}[2]{{[#1;#2[}}
\providecommand{\dotfill}{\ldots}

\providecommand{\ds}{\displaystyle}
\providecommand{\textding}[1]{#1}
\providecommand{\Oij}{\(\mathrm{O};\vec{\imath},\vec{\jmath}\)}
\providecommand{\Oijk}{\(\mathrm{O};\vec{\imath},\vec{\jmath},\vec{k}\)}
\providecommand{\Ouv}{\(\mathrm{O};\vec{u},\vec{v}\)}
\providecommand{\intervFF}[2]{\ensuremath{[#1;#2]}}
\renewcommand{\euro}{\ifmmode\text{\texteuro}\else\texteuro\fi}
\DeclareUnicodeCharacter{00B0}{\ifmmode^{\circ}\else\textdegree\fi}
\DeclareUnicodeCharacter{20AC}{\ifmmode\text{\texteuro}\else\texteuro\fi}
% Énoncés : Éducation nationale. Transcription et corrigés : APMEP, Denis Vergès et contributeurs. Découpage et classement : CoopMaths, Sébastien Lozano, Éric Elter et contributeurs. Adaptation de la présentation : Maths Explorer.
% CC BY-SA : https://coopmaths.fr/www/about/
% https://coopmaths.fr/alea/static/bac/2024/tex/bac_2024_05_sujet2_ameriquenord_1.tex
\begin{document}
\noindent\textit{APMEP / CoopMaths. CC BY-SA. Présentation adaptée par Maths Explorer.}\par\medskip
Les données publiées le 1\up{er } mars 2023 par le ministère de la transition écologique sur les immatriculations de véhicules particuliers en France en 2022 contiennent les informations suivantes :

	\begin{itemize}
		\item $22,86\,\%$ des véhicules étaient des véhicules neufs ;
		\item $8,08\,\%$ des véhicules neufs étaient des hybrides rechargeables;
		\item $1,27\,\%$ des véhicules d'occasion (c'est-à-dire qui ne sont pas neufs) étaient des hybrides rechargeables.
	\end{itemize}

	\emph{Dans tout l'exercice, les probabilités seront arrondies au dix-millième.}
	
\medskip
	
	\textbf{Partie A}

Dans cette partie, on considère un véhicule particulier immatriculé en France en 2022.

On note :

\begin{itemize}
	\item $N$ l'évènement \og le véhicule est neuf \fg{};
	\item $R$ l'évènement \og le véhicule est hybride rechargeable\fg{};
	\item $\overline{N}$ et $\overline{R}$ les évènements contraires des évènements contraires de $N$ et $R$.
\end{itemize}

\begin{enumerate}
\item Représenter la situation par un arbre pondéré.

\item Calculer la probabilité que ce véhicule soit neuf et hybride rechargeable.

\item Démontrer que la valeur arrondie au dix-millième de la probabilité que ce véhicule soit hybride rechargeable est \np{0,0283}.

\item Calculer la probabilité que ce véhicule soit neuf sachant qu'il est hybride rechargeable.
\end{enumerate}

\textbf{Partie B}

\medskip

Dans cette partie, on choisit $500$ véhicules particuliers hybrides rechargeables immatriculés en France en 2022.

Dans la suite, on admettra que la probabilité qu'un tel véhicule soit neuf est égale à $0,65$.

On assimile le choix de ces 500 véhicules à un tirage aléatoire avec remise.

On appelle $X$ la variable aléatoire représentant le nombre de véhicules neufs parmi les $500$ véhicules choisis.

\medskip

\begin{enumerate}
\item On admet que la variable aléatoire $X$ suit une loi binomiale. Préciser la valeur de ses paramètres.
\item Déterminer la probabilité qu'exactement $325$ de ces véhicules soient neufs.
\item Déterminer la probabilité $p(X \geq 325)$ puis interpréter le résultat dans le contexte de l'exercice.
\end{enumerate}

\textbf{Partie C}

\medskip

On choisit désormais $n$ véhicules particuliers hybrides rechargeables immatriculés en France en 2022, où $n$ désigne un entier naturel strictement positif.

On rappelle que la probabilité qu'un tel véhicule soit neuf est égale à $0,65$.

On assimile le choix de ces $n$ véhicules à un tirage aléatoire avec remise.

\medskip

\begin{enumerate}
\item Donner l'expression en fonction de $n$ de la probabilité $p_{n}$ que tous ces véhicules soient d'occasion.

\item On note $q_{n}$ la probabilité qu'au moins un de ces véhicules soit neuf. En résolvant une inéquation, déterminer la plus petite valeur de $n$ telle que $q_{n} \geqslant \np{0,9999}$.
\end{enumerate}

\end{document}
